\documentclass[10pt,a4paper]{amsart}
\usepackage[margin=19mm]{geometry}
\usepackage{amsmath,amssymb,mathtools}
\usepackage[scr=rsfs,cal=euler]{mathalfa}
\usepackage{xfrac}
\usepackage[unicode,hidelinks,bookmarks=false]{hyperref}
\newcommand{\ie}{i.e.\ }
\newcommand{\cf}{cf.\ }
\newcommand{\resp}{respectively\ }
\newcommand{\Subsection}{\subsection}
\providecommand{\nobreakdash}{\nobreak\hbox{-}}
\setlength{\emergencystretch}{2em}
\multlinegap0pt
\usepackage{bm}
\usepackage{bbm}
\usepackage{dsfont}
\usepackage{xspace}
\usepackage{cancel}
\usepackage{mathtools}
\usepackage{cleveref}
\usepackage{amsthm}
\makeatletter
\@ifundefined{proposition}{\newtheorem{proposition}{Proposition}}{}
\@ifundefined{lemma}{\newtheorem{lemma}[proposition]{Lemma}}{}
\makeatother
\title[Graphs, Axial Algebras and their Automorphism Groups]{Graphs, Axial Algebras and their Automorphism Groups}
\author{Hans Cuypers}
\date{Source: arXiv:2603.03822v1 (CC BY 4.0)}
\begin{document}
\maketitle

\noindent\emph{Source and licence.} Graphs, Axial Algebras and their Automorphism Groups. Version arXiv:2603.03822v1 (CC BY 4.0), distributed under the Creative Commons Attribution 4.0 International licence, as recorded in the arXiv licence metadata for this exact version and shown on the abstract page https://arxiv.org/abs/2603.03822v1. Full rights notice: licenses/arxiv-2603.03822-CC-BY-4.0.txt.\par\smallskip

\noindent\textit{Editorial scope.} Counted unit: the Lemma whose source label is \texttt{proper} in the article (source0.tex lines 675--678, source label \texttt{proper}), delivered together with the article's own complete original proof of it (source0.tex lines 680--751, one \texttt{proof} environment of 72 lines). No mathematics is written, repaired or re-derived: statement and proof are the article's own text line for line. Required context retained uncounted: casesfora (lines 611--623). Every definition, display and label of those lines is retained unchanged. Omitted from this package and not redistributed: 1--610, 624--674, 679--679, 752--1994. Nothing inside a retained excerpt is omitted or altered: every retained body is delivered exactly as the article's lines are written, so no omission receipt of any kind accompanies this package. Citations: the retained excerpts carry no citation command at all, so no bibliography entry of the article is transcribed and no reference is redistributed. Reference adaptation: none was needed and none was made; every reference inside the delivered unit resolves inside it, so no unresolved reference or citation is produced. Printed slips kept literal: no printed word, symbol, hypothesis or number of the counted statement or of its proof was altered. Layout and class: the article's own class is \texttt{amsart} with packages including amssymb, verbatim, mathabx, MnSymbol, enumitem, fontenc, babel, todonotes, hyperref, cleveref; the wrapper is the lane's pinned \texttt{amsart} editorial wrapper at 10pt on A4, which restates the theorem-like environments the retained text needs and adds the article's own macro definitions that the retained text needs (none), with extra wrapper packages (bm, bbm, dsfont, xspace, cancel, mathtools, cleveref). The delivered numbering is the unit's own. Assets: no retained span contains a figure environment, a graphics-inclusion command, a PDF-page inclusion, a table or a TikZ picture; no image file is copied into this package, the package declares no asset dependency and no image file is redistributed with it. Licence: version arXiv:2603.03822v1 of this article is distributed under the Creative Commons Attribution 4.0 International licence, as recorded in the arXiv licence metadata for this exact version and shown on the abstract page https://arxiv.org/abs/2603.03822v1. A full rights notice is delivered at licenses/arxiv-2603.03822-CC-BY-4.0.txt. Characters: every retained excerpt is pure ASCII, so the delivered engine, XeLaTeX with its default Latin Modern text font, reports no missing character.
\begin{lemma}\label{casesfora}
One of the following holds.
\begin{enumerate}
\item 
There is a $\lambda\in \mathbb{F}^*$ with $a=\lambda\displaystyle \sum_{y\in \mathrm{supp}(a)} y$.
For any two vertices $x,y$ of $\mathrm{supp}(a)$ we find $(x,y)\in E$ with label $\frac{1}{2-|\mathrm{supp}(a)|}\neq \frac{1}{2}$ and $|\mathrm{supp}(a)|>2$.
%In particular, $|\mathrm{supp}(a)|\neq 0,-2$ in $\mathbb{F}$.
\item

For any two vertices $x,y$ of $\mathrm{supp}(a)$ we find $(x,y)\in E$ with label $\frac{1}{2}$.
Moreover, $\displaystyle\sum_{y\in \mathrm{supp}(a)} \lambda_y=0$ in $\mathbb{F}$.
\end{enumerate}
\end{lemma}

\begin{lemma}\label{proper}
Suppose $\mathrm{supp}(a)$ is a proper subset of $X$.
Then the induced subgraph on $\mathrm{supp}(a)$ is an ideal subgraph.
\end{lemma}

\begin{proof}
Assume that $\mathrm{supp}(a)$ is a proper subset of $X$.
Then consider an element $x\in X$ which is not in $\mathrm{supp}(a)$, but adjacent to at least one vertex of $\mathrm{supp}(a)$, say $z$.
Then for $\beta\in \mathbb{F}$ we find
$$\begin{array}{ll}
xa-\beta a
&=\displaystyle\sum_{y\in \mathrm{supp}(a)} \lambda_y (xy-\beta y)\\
&=\displaystyle\sum_{y\in \mathrm{supp}(a),x\sim y}  \lambda_y(xy-\beta y)-\beta \sum_{y\in \mathrm{supp}(a),x\not\sim y} \lambda_y y\\
&=\displaystyle\sum_{y\in \mathrm{supp}(a),x\sim y} \lambda_y(\alpha_{x,y} x +(\alpha_{x,y}-\beta) y)-\beta \sum_{y\in \mathrm{supp}(a),x\not\sim y} \lambda_y y.\\

\end{array}$$

Setting  $\beta=0$, we find this product to be non-zero, and 
minimality of $|\mathrm{supp}(a)|$ implies that $x$ is adjacent to at least all but one vertices of $\mathrm{supp}(a)$.



Now assume that  $x$ is adjacent to all but one of the vertices of $\mathrm{supp}(a)$. Say $u\in \mathrm{supp}(a)$ is the unique vertex not adjacent to $x$.
Then, as $xa\neq 0$,  the support of $xa$ is of the same size  as the support of $a$.
In particular,  \cref{casesfora} also applies to $\mathrm{supp}(xa)$, so we find $\alpha_{x,y}=\alpha_{x,z}$ for all $y\in \mathrm{supp}(a)$ different from $u$. 
If both $\alpha_{x,z}$ and $\alpha_{z,u}$ are different from $\frac{1}{2}$, then, by \cref{casesfora} applied to both  $a$ and to $xa$, we can assume
that (up to a scalar) $a= \displaystyle\sum_{y\in \mathrm{supp}(a)} y$ and $xa=\lambda\displaystyle\sum_{y\in \mathrm{supp}(xa)} y$ for some nonzero $\lambda\in\mathbb{F}$.
But then $xa-\lambda a=\lambda (x-u)$ and $x(x-u)=x\in I$, which is against our assumption.
So, at least one of $\alpha_{x,z}$ or $\alpha_{z,u}$ equals $\frac{1}{2}$.

 
If $|\mathrm{supp}(a)|>2$, then  $\mathrm{supp}(a)$ and $\mathrm{supp}(xa)$ intersect in at least two vertices implying that all
edges inside their union have the same label $\frac{1}{2}$.
We find  $xa-\frac{1}{2} a\in I$ has support of size less than $|\mathrm{supp}(a)|$ and hence  is $0$.
But then   
$$\displaystyle\sum_{y\in \mathrm{supp}(a),y\sim x} \lambda_y =\displaystyle\sum_{y\in \mathrm{supp}(a),y\neq u} \lambda_y=0.$$

Above we found that also $\displaystyle\sum_{y\in \mathrm{supp}(a)} \lambda_y=0$
and hence $\lambda_u=0$, which is impossible.


In the case where $|\mathrm{supp}(a)|=2$ we can, without loss of generality, assume $a=z+\lambda_u u$.
But then $xa= \alpha_{x,z}(x+z)$, and $x+z\in I$.
So also $x(x+z-a)= x(x-\lambda_u u)=x\in I$, which is against our assumptions.

We conclude that $x$ is always adjacent to none or all vertices in $\mathrm{supp}(a)$. 



So, assume that $x\in X$ is not in $\mathrm{supp}(a)$, but $x$ is adjacent to some fixed $z\in\mathrm{supp}(a)$
and hence to all $y\in \mathrm{supp}(a)$.
Setting  $\beta=\alpha_{x,z}$,  in the equation
$$ xa-\beta a=\displaystyle\sum_{y\in \mathrm{supp}(a),x\sim y} \lambda_y(\alpha_{x,y} x +(\alpha_{x,y}-\beta) y),$$ 
we find $xa-\beta a\in I$ to have a support of size at most 
$|\mathrm{supp}(a)|$.

If $xa-\beta a=0$, then $\alpha_{x,y}=\beta=\alpha_{x,z}$ for all $y\in \mathrm{supp}(a)$
and $$0=xa-\beta a=(\displaystyle\sum_{y\in \mathrm{supp}(a)}\lambda_y) x\in I.$$
Hence $\displaystyle\sum_{y\in \mathrm{supp}(a)}\lambda_y=0$.
If, moreover, $\alpha_{z,z'}\neq \frac{1}{2}$ for some $z,z'\in \mathrm{supp}(a)$, then, by \cref{casesfora} we have $a_{z,z'}=\frac{1}{2-|\mathrm{supp}(a)|}$ and 
there is a $\lambda\in \mathbb{F}^*$ with $a=\lambda \displaystyle\sum_{y\in \mathrm{supp}(a)} y$.
So, $\lambda |\mathrm{supp}(a)|=0$ in $\mathbb{F}$, contradicting $\alpha_{z,z'}\neq 0, \frac{1}{2}$.
Hence $\alpha_{y,y'}=\frac{1}{2}$ for all $y,y'\in \mathrm{supp}(a)$ and we find $\mathrm{supp}(a)$ to be an ideal subgraph of $\Gamma$.


If $xa-\beta a\neq 0$, then $xa-\beta a$ is also an element of $I$ with minimal support equal to $\{x\}\cup \mathrm{supp}(a)\setminus \{z\}$.
But then, by \cref{casesfora}, we have  $\alpha_{x,y}=\alpha_{x,y'}$ for all $y,y'\in \mathrm{supp}(a)$ different from $z$. 
So, for fixed $z'$ in $\mathrm{supp}(a)$ different from $z$ we find that $\alpha_{x,z'}\neq \alpha_{x,z}$ and   $xa-\alpha_{x,z'}a\neq 0$ has support contained in $\{x,z\}$.
As $xa-\alpha_{x,z'}a\neq 0$, we find that $\mathrm{supp}(xa-\alpha_{x,z'}a)=\{x,z\}$ has size $2$. But then $\mathrm{supp}(a)=\{z,z'\}$ and 
$\mathrm{supp}(xa-\beta a)=\{x,z\}$.

As $xa-\beta a$ as well as $xa-\alpha_{x,z'}a$ are in $I$ and have support of size $2$, we find, by \cref{casesfora}, that the edges $(x,z)$ and $(x,z')$ both have label $\frac{1}{2}$, contradicting that they have different label.

If there is no $x\in X$  which is not in $\mathrm{supp}(a)$ but  adjacent to an element from $\mathrm{supp}(x)$, then, by the weakly connectedness
of $\Gamma$ there is an $x\in X$ but not in $\mathrm{supp}(a)$ for which there is a $z\in \mathrm{supp}(a)$ adjacent to $x$.
Considering products of the form $ax-\beta a$ we obtain the same conclusions as above.
\end{proof}

\end{document}
